\problema{Word Ladder}{127}{src/01-grafos/127-word-ladder.cpp}{H}

Dadas \texttt{beginWord}, \texttt{endWord} y una lista \texttt{wordList},
buscamos la longitud de la secuencia de transformación \emph{más corta} de
\texttt{beginWord} a \texttt{endWord}, donde cada paso cambia exactamente una
letra y cada palabra intermedia debe existir en \texttt{wordList}. La
longitud cuenta las palabras (incluyendo begin y end); si no existe tal
secuencia, o \texttt{endWord} no está en \texttt{wordList}, devolvemos
\texttt{0}.

\paragraph{La idea, con un ejemplo de la vida real.} Imagina un juego de mesa
donde cada palabra es una casilla. Dos casillas están conectadas por un puente
si sus palabras se diferencian en una sola letra (por ejemplo \texttt{hit} y
\texttt{hot}). Queremos ir de la casilla inicial a la casilla final dando la
menor cantidad de saltos posible.

Para encontrar el camino más corto en un tablero así, la mejor forma es
explorar ``en anillos'', desde la casilla de salida hacia afuera: primero todas
las casillas a las que llegas en un salto, luego todas a las que llegas en dos
saltos, y así. En cuanto uno de esos anillos toca la casilla final, sabemos que
ese es el número mínimo de saltos, porque exploramos de menos a más. Esa forma
de explorar por anillos, o por capas, se llama \emph{BFS}:
cada anillo del BFS es un paso de transformación.

\begin{center}
\ttfamily
\begin{tabular}{c}
hit \\
$|$ \\
hot \\
\end{tabular}
\quad
\begin{tabular}{cc}
dot & lot \\
$|$ & $|$ \\
dog & log \\
\end{tabular}
\quad
\begin{tabular}{c}
cog \\
\end{tabular}
\end{center}

\paragraph{Cómo lo resuelve el código, paso a paso.} No dibujamos todos los
puentes de antemano, porque con muchas palabras sería carísimo. En vez de eso,
descubrimos los vecinos de una palabra ``sobre la marcha''.
\begin{enumerate}
  \item Usamos una \emph{cola} (una fila: primero en entrar, primero en salir)
        y metemos la palabra inicial.
  \item Para sacar los vecinos de una palabra, recorremos cada posición y
        probamos cambiarla por las $26$ letras del abecedario. Cada palabra que
        resulte y que esté en \texttt{wordList} es una casilla vecina. Guardamos
        las palabras en un \texttt{unordered\_set} (una especie de bolsa mágica
        donde revisar si algo está adentro es instantáneo).
  \item Vamos anillo por anillo: al empezar cada nivel miramos cuántas palabras
        hay en la cola en ese momento (\texttt{cola.size()}) y procesamos justo
        esas antes de pasar al siguiente anillo.
  \item Marcamos como visitada cada palabra que ya usamos, para no volver a
        pasar por ella y dar vueltas en círculo.
\end{enumerate}
En cuanto un anillo alcanza \texttt{endWord}, contamos cuántos anillos llevamos
y esa es la respuesta.

\lstinputlisting{src/01-grafos/127-word-ladder.cpp}

\complejidad{$O(N \cdot L^2 \cdot 26)$}{$O(N \cdot L)$}

\paragraph{¿Qué significa esa complejidad?} Aquí $N$ es cuántas palabras hay en
\texttt{wordList} y $L$ cuántas letras tiene cada palabra. Para cada palabra
probamos cambiar cada una de sus $L$ letras por las $26$ del abecedario, y armar
cada palabra nueva cuesta otras $L$ letras; por eso aparece $L^2$ y el $26$.
En resumen: es el trabajo de generar y revisar todos los cambios de una letra,
por todas las palabras.

\paragraph{Consejos para la entrevista (Amazon).} Lo primero, sin falta, es
revisar que \texttt{endWord} esté en \texttt{wordList}; si no está, la respuesta
es \texttt{0} y ni siquiera vale la pena empezar a explorar. El error típico es
buscar los vecinos comparando la palabra contra cada palabra de la lista letra
por letra, lo cual es lentísimo cuando la lista es grande; es mucho mejor
generar los cambios de una letra y preguntarle a la bolsa mágica
(\texttt{unordered\_set}) si esa palabra existe. También puedes mencionar el
\emph{BFS bidireccional}: explorar al mismo tiempo desde la palabra inicial y
desde la final, para que los dos anillos se encuentren en medio y haya menos
trabajo.
